% =====================================================================
%  R. Nielsen, A FEW WORDS OCCASIONED BY PROF. SCHARLING'S APOLOGY FOR DR. MARTENSEN'S DOGMATICS.
%  Copenhagen: C. A. Reitzel, 1850.  ENGLISH TRANSLATION, from transcription.tex (Danish) as rebuilt
%  2026-10-06 (see its header).
%  GENERATED by ebook-method/scharling-apologie/tr_build.py --write from tr_template.tex and the
%  translator's part in texts/nielsen/scharling-apologie/.parts/tr/ (TRANSLATOR-BRIEF.md).
%  Conventions: TRANSLATION-PLAYBOOK.md.  Page numbers of the 1850 print in the margin.
% =====================================================================
\documentclass[12pt,a4paper,openany]{book}
\usepackage[utf8]{inputenc}
\usepackage[T1]{fontenc}
\usepackage{libertinus}
\usepackage{libertinust1math}
\usepackage{textalpha}
\usepackage[english]{babel}
\usepackage{enumitem}
\usepackage[protrusion=true,expansion=true]{microtype}
\usepackage{setspace}
\usepackage[margin=1.3in]{geometry}
\usepackage{fancyhdr}
\usepackage{hyperref}
\hypersetup{pdftitle={A Few Words Occasioned by Prof. Scharling's Apology for Dr. Martensen's Dogmatics}, pdfauthor={Rasmus Nielsen}, hidelinks}
\setlength{\headheight}{15pt}
\pagestyle{fancy}
\fancyhf{}
\fancyhead[LE,RO]{\thepage}
\fancyhead[C]{\textit{R. Nielsen: A Few Words on Prof. Scharling's Apology (1850)}}
\renewcommand{\thefootnote}{\ifcase\value{footnote}\or *)\or **)\or ***)\else
  \arabic{footnote})\fi}
\setstretch{1.3}
\usepackage{marginnote}
\newcommand{\opage}[1]{\marginnote{\footnotesize\textit{[#1]}}}

\begin{document}

\begin{titlepage}
\begin{center}
\vspace*{3em}
{\Large A Few Words}\\[1em]
occasioned by\\[1em]
{\Huge Prof. Scharling's Apology}\\[1em]
for\\[1em]
{\LARGE Dr. Martensen's Dogmatics}\\[1em]
by\\[1em]
{\large Lic. R. \emph{Nielsen,}}\\
{\small Professor of Philosophy.}\\[2em]
\rule{8em}{0.4pt}\\[2em]
{\large Copenhagen.}\\[0.5em]
University Bookseller C. A. \emph{Reitzel.}\\[0.5em]
{\small Printed by Royal Court Printer Bianco Luno.}\\[0.5em]
1850.\\[4em]
{\small Translated from the Danish}
\end{center}
\end{titlepage}
\thispagestyle{empty}
\vspace*{12em}
\begin{center}
\begin{minipage}{0.8\textwidth}
\textit{Μιλήσιοι ἀξύνετοι μέν οὐκ εἰσί· δρῶσι δὲ, οἷαπερ οἱ ἀξύνετοι.}\\[0.4em]
\hspace*{\fill}\textbf{Demodacus.}\\[1em]
{\small [The Milesians are not stupid, but they act just as the stupid do.]}
\end{minipage}
\end{center}
% the motto is the epigram of Demodocus (the print: „Demodacus“); the bracketed English is the translator's
\clearpage

% --- p. 3 ---
\setcounter{footnote}{0}%
\opage{3}In ``Nyt theologisk Tidsskrift'' (vol. 1, no. 2, pp.
348---375) His Reverence Dr. theol. Prof. C.
E. Scharling has --- remarkably enough! --- appeared as
apologist for Dr. Martensen's dogmatics.

That so energetic a step on Prof. Scharling's
part has come somewhat unexpectedly to me and has also put me
in great embarrassment, I will not conceal. For a
sharp scientific rejoinder from Dr. Martensen himself
I have long been prepared, but for an Apology from Prof.
Scharling, no, that I confess, for so genial an
Apology I was not prepared at all.

Why I call the Scharlingian Apology \emph{genial}
is easy to explain. A scientific defense
holds strictly to the opponent's objections and grounds;
but such formal considerations are naturally of subordinate
importance in an Apology whose entire probative force can
be reduced to the categories: goodwill and ill will, approval
and disapproval, praise and blame.

Where Prof. Scharling at present has his scientific
standpoint, and whether he has one at all, I shall
not be able to decide; on the other hand, from the present
% --- p. 4 ---
\setcounter{footnote}{0}%
\opage{4}document it is by no means difficult to state his apologetic
standpoint: Prof. Scharling is at the present time ---
genially praising Prof. Martensen, genially censuring
his opponents. What Prof. Scharling now actually
\emph{thinks} about the relation of ``faith'' to ``knowledge'', about the relation of religion
to philosophy, of Christianity to theology,
the little acuteness that was granted me
will not be able to discover from his genial Apology;
what I can discover without effort, on the other hand, and what
the theological professor has undeniably also succeeded
in presenting with rare theological clarity, is the
turning point, surely by no means unimportant for ``theological science'' either,
that His Reverence has become angry
with Mr. Stilling and me.

Thus the first faculty-sanctioned yield of
the first faculty-sanctioned Apology for our newest literature's
first faculty-sanctioned dogmatics is to be expressed and
held fast in the proposition, as weighty as it is momentous: \emph{Professor Scharling has become angry with Mr.}
\emph{Stilling and Professor Nielsen}.

That on this serious and solemn occasion I
name Mr. Stilling first has its ground
in modesty, and is done solely for politeness' sake; for
according to the unequal, though surely fair, distribution of the anger,
the first place properly belongs to me.
The matter is this: with Mr. Stilling Professor
Scharling is angry for one reason only, namely that Mr.
Stilling has permitted himself to write against Dr.
Martensen; with me, on the other hand, His Reverence is angry
for two reasons, for I have not only been so bold
as to write against Dr. Martensen, but had besides previously
% --- p. 5 ---
\setcounter{footnote}{0}%
\opage{5}been so foolhardy as to write a good deal against ``the theological
mediocrities.''\footnote[1]{Cf. ``The Gospel-Faith and the Modern Consciousness''.}

There is, then, in the genial Apology a common chastisement
that befalls us both. We have both grievously
sinned against theological science by ``the unpleasant
surprise'' that our ``more detailed assessments'' of
``the Christian Dogmatics'' have prepared for Dr. Martensen;
``the scornful irony, the teasing mockery with which the verdict of rejection
is accompanied, the cruel pleasure with which
we'' --- that I too may express myself genially --- ``shoot off the arrows of irony'' at the poor ``celebrated'' author,
``seem (p. 351) unable to exclude the supposition
that motives other than the pure zeal for truth appealed to
have made themselves felt in the composition of our
critiques''; in short: we are both chastised, and
thoroughly, for we have both been naughty.

But besides the common theological chastisement thus
awarded to us both, Prof. Scharling's apologetic geniality has yet another chastisement, a surplus, an
extra helping of chastisement, which is awarded to me in particular. Although we
are, as said, both naughty, Mr. Stilling is nevertheless
not nearly so naughty as I. Admittedly Mr.
Stilling is violent --- ``perhaps even the most violent'' ---
but notwithstanding (cf. Anm. p. 360) ``Mr.
Stilling's writing nevertheless bears, alongside personal touchiness,
the stamp of a subjective conviction that awakens the reader
to seriousness and sympathy''\dots{} It is otherwise with
me! ``It is otherwise with Prof. Nielsen. His
presentation calls forth the impression as if he might well
% --- p. 6 ---
\setcounter{footnote}{0}%
\opage{6}make the offer: ``If you will now take my opinion, then
I will in turn defend your opinion''\dots{} --- Alas, what shall
I answer? It would perhaps be better if I did not answer at all,
but kept silent, and bowed in all deference
to the theological chastisement! for, seriously speaking, what shall
I answer? With me things are going to the extreme; to my
astonishment and most unpleasant surprise I now learn
that I stand at the edge of two extremes,
two dreadful extremes, whose names I dare not even
mention, but must leave the reader to discover from Prof.
Scharling's description.

The first extreme is described (p. 364 Anm.) thus:
``He (Prof. N.) does not readily name the theologians
without some scornful epithet or other (school theologians,
middle-way theologians, \&c.); everything stupid and wrong
that theologians may have deserved and not deserved,
he gathers into features of the caricatures with which he seeks
to make his theological opponents martyrs of ridicule.
This, by the way, is connected with a strange
illusion about the esteem which he thinks
theologians, especially theological doctors and professors, enjoy.
Therefore he persecutes these with the same bitterness and
satire as Voltaire (sic!) in his time the priests.''

The second, still far more questionable extreme
is described thus: after Mr. Stilling's proposition
\textit{quamvis absurdum sit, credo} has been touched on provisionally, it says
in the text (p. 364): ``nevertheless there must
underlie the latter declaration, if it is not to be dismissed
out of hand as excluding all further rational
deliberation, the notion that the Christian dogmas,
however unreasonable they may present themselves to the subject,
nevertheless have an objectively valid and eternal
% --- p. 7 ---
\setcounter{footnote}{0}%
\opage{7}truth, which it is only not possible, on account of hindrances
or limitations in the subject'' (---\textit{NB}---)
``to cognize in a given state.'' Upon this psychological
explanation there now follows the note, and the note one must
by no means skip, for the note is the best part, for
the note is even more genial than the text, for in
the note lies the point, for in the note Professor
Scharling says: ``Whether this (--- namely what is adduced
about Mr. Stilling's proposition ---) should, however,
be Prof. Nielsen's opinion, I do not know; for he finds
Stilling's cited formula still too ambiguous, and wants
it interpreted in the proposition: \textit{credo, quia absurdum est.}
Thus there might here have come about in actuality what
J. P. Lange (Christl. Dogmatik, 1849, 1ster Th. p.
409) has recently posited as a \emph{possibility}, ``daß am
Ende auch noch ein Philosoph kommen könnte, der das
Denken oder das Bewustseyn selbst ebenfalls für eine Illusion
erklären wollte. Was wäre nun da anzufangen?
Man mußte ihn mit dem Zugeständniß, daß sein eignes
Bewustseyn in Taumel gerathen sein könne, nach Hause
schicken.''

Nach Hause! so then: nach Hause! yes, let it rest
there: nach Hause!

When with my ``Gospel-Faith'' and my ``examining
review'' I advanced against a theological science that is surely
very considerable in the literature, but in my conviction
wholly untenable, I was even then not
prepared for so few possibilities; but that this \emph{possibility},
this here \emph{possibility} of which Prof. Scharling
speaks, should be just so near to ``having come about
in actuality'', no, that I confess, for that I was
% --- p. 8 ---
\setcounter{footnote}{0}%
\opage{8}not at all prepared: that is indeed also a ``surprising'' stratagem
of the genial Apology.

At this stratagem I and my writings must then
feel roughly like the council of war in a middling
fortress that is unexpectedly taken by surprise by a strong
and cunning enemy; one foresees clearly that resistance is
futile, and yet, since it is after all not an open
town, but in all ``mediocrity'' a fortress, one feels that
one must, at least for honor's sake, do something. So,
yes precisely so, it now goes with me in this
apologetic surprise attack. Prof. Scharling throws down the gauntlet
to believing speculation: I feel that I must
do something, and foresee nevertheless that nothing can come
of it.

That I now really, despite all my ``irritability'', must have
a certain reason to be cautious with Prof. Scharling's
literary verdicts, His Reverence himself will be able to testify
when he thinks back. Some years ago, it
was just at the time when --- as His Rev. himself says
p. 352 --- ``the so-called speculative theology here in this country
enjoyed an esteem whose luster had already begun
to fade in its original home,'' I published a
commentary on Paul's Epistle to the Romans, composed not in a
strictly Hegelian but in a speculative-theological spirit. This commentary
at once had the misfortune to attract Prof. Scharling's
attention, for which reason he then also, after his genial
custom, found himself entitled to deliver a literary
verdict, so as at the same time to set me a little
to rights.

Had I then felt a calling to engage
with Prof. Scharling, the occasion would have been as favorable
as possible. ``In those days'' --- says His Rev. himself p.
% --- p. 9 ---
\setcounter{footnote}{0}%
\opage{9}353 --- ``a harsh judgment generally fell upon all men of science
who would not or could not put themselves in
immediate relation to the prevailing system''; so then:
in those days, when the reasoners' objection to speculative
cognition was ``dismissed with a stern rebuke
for spiritual narrowness and unfamiliarity with
the essence of philosophy'', in those days it was precisely that Prof.
Scharling dispatched my attempt on the Epistle to the Romans\footnote[1]{Tidsskrift for Litteratur og Kritik, ed. F. C. Petersen. 1842. vol. 1. p. 210.},
and in those days I would then also, given my exceedingly
great ``irritability'', have had the finest opportunity to
be --- a little naughty toward Prof. Scharling.

Why then, since my ``irritability'' is after all
so excessive, did I neglect the favorable opportunity, and
bow in modest silence under the genial chastisement? --- Dear reader, what shall I answer? If I answer
that I did it because I felt a certain diffidence about
attacking a man who shortly before had been my teacher:
then I will hardly be believed; for His Reverence
already at that time accused me of ironizing presumption.
If I answer that I did it because I preferred to
spare a man who evidently by a pure mistake
had come to deliver his verdict in a matter that
concerned theological speculation: then I will hardly be believed
either; for I am after all once and for all suspected of ``other
motives''. Well, then I will answer what is true, and
give an answer that Prof. Scharling will surely believe better
than any other; I will then answer him and say it straight
out: I did it out of fear.

% --- p. 10 ---
\setcounter{footnote}{0}%
\opage{10}What instilled this fear in me was by no means
the sharpness of the Scharlingian objections (for they lacked
all sharpness), nor the weight of the Scharlingian grounds
(for they lacked all weight), but a fatal circumstance,
namely that Prof. Scharling in my youthful attempt
could find no --- ``spirit'' at all. Yes,
it was precisely this circumstance that already at that time
paralyzed my strength, and struck into me a terror I have not
since that time got over.

For a more thorough appreciation of my great embarrassment
I can perhaps best lead the reader by a
small detour.

Before one engages in a literary dispute at all,
one should sensibly look about to see whom one is going
to deal with, and, as far as possible, seek to sound out the polemical
forces. As regards scientific polemic in particular,
it seems to me as if authors could best
be divided into the following three classes: 1) the authors
who, in the form of genius and of freedom, are sure of
principle and idea; 2) the authors who recognize the principle
in the form of a fundamental tenet, and therefore must bow
to regular consistency; 3) the authors who do not
step aside for any principle, since in their literary
geniality they do not really know what a principle is.

That authors of the first rank are dangerous polemicists,
especially when they are also gifted with so witty
an imagination that they can even hide the gravity of the serious
in the lightness of jest, goes without saying. I
name here Mr. Kierkegaard and Prof. Heiberg, for by
naming these names I relieve the reader and myself
of a prosaic description of the higher polemic.

To polemicists of the second class belong all capable
% --- p. 11 ---
\setcounter{footnote}{0}%
\opage{11}and reliable men of science, who seek their strength directly % print: paalideliae (misprint)
in proofs and grounds; these authors are as a rule
bound to the prosaic form of formal seriousness,
and are therefore also best known under the name of ``the serious
authors''.

To the polemicists of No. 3 belong all principle-free,
or rather: unprincipled, yet serious, indefatigable and
learned writers, all worthy members of the worthy
family \emph{Wagner}\footnote[1]{Cf. Goethe's Faust.}, who in so worthy a manner % print: Gocthes (misprint)
advance science in all directions and adorn literature
in all countries. When an author of this kinship
makes as if to engage in a literary dispute, especially in a
dialectical dispute over principle, one must be careful, very careful,
for this kind of folk is extremely dangerous. --- With the polemicists
of the first rank one still gets off reasonably well,
provided one otherwise has a good cause and a sound
thought (even if one loses, the loss is always a gain); but
woe to him who must defend himself with a principle against those
who are not hit by any principle; woe to him who must defend
the unity of a thought against those who appropriate all manner of thoughts
and all sorts of ideas, and precisely thereby
prove in the finest way ``that the spiritual nature of human beings
constitutes a unity''; yes, woe to him who falls into the hands
of a polemicist of No. 3, for --- with the Wagner
family there is no getting along.

The reader surely already senses my embarrassment. --- Under which rubric Prof. S. will place me is easy to see
from the ominous remark: ``nach Hause,'' for me, on the other hand,
it is by no means easy to see under which rubric I am in this
% --- p. 12 ---
\setcounter{footnote}{0}%
\opage{12}matter to place Professor Scharling. To put him
among the polemicists of the first rank will hardly
do (now of course, if someone else wants to do it,
I shall say nothing against it, but --- I dare not).
To place His Reverence under No. 3 will do still
less; that would indeed be a bloody insult to
a man who is not only a member of the theological faculty,
but, if I remember rightly, even a member of
the Society of Sciences. My serious endeavor must
therefore aim at getting him placed in a class with the
``serious men of science'' under rubric No. 2; but
how thankless this endeavor is, and how little
hope I can have this time of bringing it to any
favorable result, the expert will easily be able to see
from a cursory reading of his genial Apology.

As one of the main features that especially characterize
this Apology, and by which Prof. Scharling together with
his theological journal seems especially to find himself in --- ``a
steady progress'', I must here emphasize the literary
unceremoniousness, an interesting main feature that is likewise connected
with the genial.

``It is strange,'' says Prof. Scharling, as he
(p. 355) mentions Mr. Kierkegaard, ``that no one has yet
made this author's writings the object of a
more detailed analysis, which they both deserve in themselves, and
for which there seems to be special occasion in the approval
they have won among readers
of the most different levels of education.'' Since I am now wholly in agreement
with Prof. Scharling in this, it seems to me for precisely that reason
most ``strange'' that His Rev. has been able with so much unceremoniousness
to engage in so rather dismissive an
% --- p. 13 ---
\setcounter{footnote}{0}%
\opage{13}assessment of these writings, without there being in this
assessment the slightest trace that His Reverence % print: Høiærvædighed (misprint)
himself has deigned to give them even a cursory, let alone
a more detailed, analysis.

That Prof. Scharling has not deigned to give Mr.
Kierkegaard's books any analysis at all can be seen at once from the
very first lines, in which he begins to review their
form. There is hardly any bright student who does not
already know that the Kierkegaardian writings, as regards form,
divide into two main classes, and that each writing
in each of these classes again has a form that is, in its kind, infinitely
individualized, carefully articulated and nuanced down to the finest
details. This difference of form,
however, does not bother Prof. Scharling. Without
troubling himself about some otherwise unusual transformations,
His Rev. without ceremony lumps all the books into
one heap, and then continues in genial unceremoniousness to
speak of ``the formless prolixity, the irony that jests both in season
and out of season, the forced hunt for paradoxes
and antitheses, and the loquacity, witty perhaps, yet often tiring,
that seems to prevent the author, despite
his astonishing productivity, from bringing his reflections
to any satisfactory conclusion.'' With
this conclusive conception of the Kierkegaardian form
(a conception one has probably heard before from ``busy
people'' in town) it is then also no doubt connected
that Mr. Kierkegaard's dialectic --- strangely enough! ---
in Prof. Scharling (p. 358) ``leaves a feeling
as if his presentation were not really in earnest, but only
an exercise in the higher gymnastics of the spirit.''\footnote[1]{A ``more detailed analysis'' of one of the pseudonyms (e.g. of Johannes Climacus) would, if it could only succeed reasonably well, certainly furnish strong contributions to answering the question: \emph{What is seriousness}?} % print: t. Ex. (misprint)

% --- p. 14 ---
\setcounter{footnote}{0}%
\opage{14}However, about form, and in particular about inwardness's indirect
form of communication in doubly reflected categories,
I, who myself give such hard offense against form, am unworthy
to take issue with Prof. Scharling, whose tasteful
journal not only tastes form, but also tastes
of form; a more detailed discussion of form would besides
be out of place here, and lead us too far away from
apologetic geniality. On the other hand I would --- even if
on this occasion I should give a little offense against form
--- yet most heartily like to give Prof. Scharling the well-meant
advice never again to occupy himself with
evaluating or criticizing the things he has not yet properly
understood, and to that end draw his attention to some
points in which his literary unceremoniousness seems to be
of such a nature that it cannot even be excused by
geniality.

When Prof. Scharling (p. 368) evaluates Mr.
Kierkegaard's polemic against ``the system'', he expresses himself
in such a way that anyone not otherwise more closely
informed must necessarily assume that Mr. Kierkegaard
gradually strayed from the justified standpoint
with which he began, and ended in an evident
one-sidedness. Thus it says: ``So long as Mr.
Kierkegaard's polemic confined itself to disputing
that the system was capable of taking the place of faith\dots{} --- then it was within its rights, and one followed it with
approval and interest. But the more it aimed at
removing faith and knowledge from each other, the more it was led
% --- p. 15 ---
\setcounter{footnote}{0}%
\opage{15}over into the same one-sidedness that it so vigorously combated\dots{} Therefore the continued polemic too, despite
many a brilliant flash of thought, could not properly hit its object,
the more it left the more abstract regions and
entered into the given historical circumstances.''

See, thus Mr. Kierkegaard's ``astonishing
productivity'' is now dispatched at one stroke by Prof. Scharling.
What impression must such an assessment now
make? --- On a wholly naive and credulous reader
the assessment must presumably make the impression that Mr. Kierkegaard
is here placed before a higher tribunal, that the theological-scientific
superiority, which has long kept a
watchful eye on his manifold ``dialectical turns'',
now finally, when it seems to be going too far, at last
breaks its silence and directs the restless ``genius'' back
within the proper bounds. On me, however, the assessment
makes a different impression, for on me it makes the
impression as if the whole thing might perhaps be loose talk.
The paradox of faith put forward by ``Johannes Climacus''
is after all essentially the very same paradox as the one that breathes
in ``\emph{Fear and Trembling}''. That the mood, the form and
the development must be different is self-evident, since
Johannes \textit{de silentio} expectorates in ``\emph{dialectical}
\emph{lyric}'', whereas Johannes Climacus on the other hand delivers an
``\emph{existential contribution}''.

I state this assertion without circumlocution, not only
because I think I am sure of my case, but especially
because it is really Prof. Scharling on whom,
as regards the denial, one may rightly lay the \textit{onus}
\textit{probandi}. If, then, Prof. Scharling will stand by his word,
and can clearly prove to me that Mr. K. has in the
said writings changed standpoint, that the paradox % print: Standpunkt. (misprint)
% --- p. 16 ---
\setcounter{footnote}{0}%
\opage{16}posed by Johannes \textit{de silentio} is an entirely different one and
in an entirely different relation to the \textit{Absurdum} than the one % print: Forhold. (misprint)
posed by Johannes Climacus\footnote[1]{For that Johannes \textit{de silentio} speaks of Abraham's faith, Johannes Climacus on the other hand of Christian faith, makes no difference here, where the point is precisely to demonstrate the sphere of faith.}; then, yes then, I shall not
only be willing in all geniality to let someone ``nach
Hause schicken'' me, but even into the bargain make amends to the ``Apologetics''
and gild ``the critical introduction''.

Since, however, it might easily happen that Prof. S.
would remain in debt for the proof, and precisely thereby come to prove a little
more than he actually wanted, I shall, for the sake of this case,
not omit here to adduce something which at least
in my eyes speaks considerably in his excuse. When
His Rev. read through the ``more detailed assessments'' of
Dr. Martensen's dogmatics, it no doubt happened to
him that on behalf of ``dogmatics'' as well as
of ``theology'' as a whole he felt most unpleasantly % print: Tyeologiens (misprint)
affected. In his worried mood he has then by an
entirely natural association of ideas come to think
first of Mr. Kierkegaard's paradox, then of my paradox,
then of Mr. Stilling's paradox, thereafter by a
likewise entirely natural association of ideas of ``Proudhon'',
``Max Stirner'', yes, of Proudhon's and Max % print: Stirmer (misprint)
Stirner's ``consorts'' together with the remaining ``vain % print: Stimers (misprint)
mediocrities'' (\textit{les mediocrités ambitieuses}), who
``increase the confusion by each casting his paradox out into
the world''; in so worried a moment it has then seemed to
him as if he saw Pandora's open box and saw distinctly
that the box was full of --- nothing but paradoxes;
% --- p. 17 ---
\setcounter{footnote}{0}%
\opage{17}nor in so worried a moment was it
easy for His Rev. to make sense of the paradoxes, but all paradoxes,
Christian, un-Christian, socialist, communist,
Danish, French, German, Spanish, in short: all the world's real
and unreal paradoxes have shown themselves to him in
such multiplicity and have strayed so among
one another that in the end it was not at all possible for him
to tell which was which.

A main feature that likewise characterizes the Scharlingian
Apology is a certain literary forgetfulness, a genial
forgetfulness, which in several points betrays that
the Right Reverend author surely cannot always have
the most faithful memory. Why Prof. Scharling is not
exactly any great patron of philosophical
speculation, he himself confesses, and I thus need
not explain the reason to him further; but that Prof.
Scharling (if he has not otherwise quietly changed
his earlier opinion) is not any particular
patron of theological speculation either, I could even
prove to him rather easily, if only I dared
rely a little on his literary memory.

The view of the relation between faith and knowledge
that underlies Dr. Martensen's dogmatics is
certainly the same view as the one that underlies
my exposition of Paul's Epistle to the Romans\footnote[1]{By this, however, I by no means wish to have asserted that my form of exposition should be as classical as Dr. Martensen's (I am after all no ``Genius''), nor in that youthful attempt was I able to work together quite so many heterogeneous elements as Dr. Martensen has done with so much plastic power in the finished work, in which ``a mystical breath of life exercises a secretly captivating power over the unbiased reader''. What I assert is only this, that the ``fundamental view'' is essentially the same, and I abide by this assertion of mine until I see a little proof of the opposite.}.

% --- p. 18 ---
\setcounter{footnote}{0}%
\opage{18}That Prof. Scharling must, however, have forgotten this
little circumstance is clear enough from the words with which he
(p. 371) meets the objections of criticism to the
great proposition that the dogmatician is ``\emph{at once the organ of his}
\emph{science and of his church.}'' Thus it
says: ``But every such view that recognizes in
historical development something relative and mobile
is rejected by a criticism which, while it reproaches
others for having remained standing at the
standpoint once taken, has not itself broken loose
from its former standpoint, in the respect that it has
retained from it just as much the abstract conception of
``absolute knowledge'', ``objective certainty'', ``coherent cognition'', as the notion of Christian faith
either as a closed sum of certain unalterable
doctrines or as a wholly unclear and indeterminable
feeling''. --- If Prof. Scharling, when he wrote down
these remarkable lines, had only stopped to recollect a little,
and called to mind the concept of ``absolute knowledge'',
``objective certainty'', ``coherent cognition'' that I have, e.g.,
developed in my propaedeutic logic, \textit{item} the theological
principle of cognition that runs through my Epistle to the Romans,
His Rev. would certainly have struck out the whole
passage, as not fitting me at all, although it
is precisely I who attacked that great proposition
of Dr. Martensen's. For should it be true, what
% --- p. 19 ---
\setcounter{footnote}{0}%
\opage{19}Prof. Scharling asserts, that I now have a wholly abstract
conception of ``absolute knowledge'', ``objective certainty'', % print: objectiv Vished“ (misprint)
``coherent cognition'', then it must at
least be untrue, what Prof. Scharling nevertheless in the same
breath likewise asserts, that in this respect I
have not been able to break loose from my former standpoint,
for in the two writings cited my conception of ``absolute
knowledge'' and ``objective certainty'' is after all not so entirely
abstract; but there is, once and for all, a fatality about
literary forgetfulness!

When Prof. Scharling reviewed my exposition
of the Epistle to the Romans, the spirit in theological speculation
was to him still a foreign and unknown quantity,
and therefore he --- to my astonishment ---
could not find the spirit either. Now, on the other hand, when Dr.
Martensen ``has with his usual clarity expressed himself in the introduction
to his dogmatics in a manner which on the
whole will be found as appealing as it is satisfying by
anyone who has not given up the conviction that the
spiritual nature of man constitutes a unity'', now
Prof. Scharling suddenly catches sight of the spirit, acknowledges
the principle, and then appears in all geniality as apologist
for theological speculation.

But how shall I now explain this strange phenomenon to myself?
I by no means assume that Prof. Scharling intentionally
wanted to be partial or to do me wrong; rather
I believe I have discovered a point of view from which
the whole thing can be explained in a manner as natural as it is
genial. The point lies in what His Rev. calls
``the constellation of spiritual movements''. As little as Prof.
Scharling can ever be assumed to have been partial
against ``speculation'', so little can he yet be assumed
% --- p. 20 ---
\setcounter{footnote}{0}%
\opage{20}to be partial against ``the paradox''. Prof. Scharling surely
at bottom likes ``the paradox'' just as much as
``speculation'', and only wishes that they would nicely know how to value
that, and be prettily reconciled with each other. When Prof. Scharling,
therefore, at present speaks a little harshly to ``the paradox''
and expresses himself a little differently about ``speculation'' than
he formerly did, the cause of this is hardly to be
sought in any essential change in Prof. Scharling's
``fundamental view''; but ``here a more natural
and worthier explanation at once offers itself, if we go back in time,
and attend to the difference in the constellation of
spiritual movements.''

In those days, when the speculation of faith was so ``haughty'',
and looked down so superciliously on the ``critical introduction'',
in those days when there generally fell ``a harsh
judgment on all men of science and writers who would not
or could not'' --- occupy themselves with speculation,
in those days Prof. Scharling was but little pleased with
the whole business of speculation, whether for that matter it
was philosophical or theological. Not without a certain inner
satisfaction His Rev. then greeted the days when a polemic as
justified as it was irresistible rose up against
the troublesome system, and made ``speculation'' comic;
to that extent Prof. S. then followed ``the paradox'' with ``approval and
interest'' --- But then came in turn the days when ``speculation''
crept to the cross, and almost on the verge of tears over
``the arbitrariness of the paradox'' poured out its grief into the
feeling bosom of ``the critical introduction''; and those days are
presumably more or less these days, and in these days Prof.
Scharling has then also felt that it is finally going too far, and
for that reason resolved to say an encouraging word to the
% --- p. 21 ---
\setcounter{footnote}{0}%
\opage{21}deeply bowed ``speculation'' and keep the
naughty ``paradox'' in check.

By no means, therefore, should one accuse Prof. Scharling
of the presumption of wanting to level the differences in the conflicting
views or, as it is called, to ``mediate'' speculation
with the paradox; no, Prof. S. does not aim so
high; he only thinks like this: each of the potencies can
be good for itself, provided only that in mutual recognition
they do not bother each other or hinder each other's ``fame''.
To Prof. Scharling's fame, then, let it be said that
one should, according to his genial Apology, by no means
take strict issue with opposite conceptions and heterogeneous
principles, but on the contrary (p. 373) ``one should contribute
to the recognition and appreciation of every endeavor
which, with enthusiasm for its goal and with modest acknowledgment
of the powers granted for its attainment,
aims at placing the truths on which
the eternal and temporal welfare of individual human beings and of the whole human race
depends in a clearer light, at grounding them
more deeply, and in this way contributing to awaken and
preserve reverence and devotion for their beneficent
influence.''

That a fine feeling can also really underlie
so genial an admonition must be admitted;
these are childlike notions; but the difficulty is that
such childlike notions prove nothing in a scientific
discussion.

If, however, I asserted that nowhere in the whole genial
Apology was any attempt made to
clarify the scientific point of dispute, I would nevertheless go
too far. There is really one place where His Rev. bethinks himself
that he really ought also to state the
% --- p. 22 ---
\setcounter{footnote}{0}%
\opage{22}problem; this place is found p. 367, and reads thus:
``It is clear that the lever in the dialectic employed
here is the straining of the concept: ``\emph{knowledge,}'' ``\emph{certainty,}''
``\emph{coherent cognition}'', to such a
degree of ``absoluteness'' and ``objectivity'' that it either
finds no application'' etc.

So in the unreasonably exaggerated straining of the concept:
``knowledge'', ``certainty'', ``coherent cognition'',
Prof. Scharling would see the ``lever'' with which
an all too obtrusive dialectic has tried to lift
``the Christian Dogmatics'' off its speculative hinges.
How it has come about that Prof. Scharling has caught sight
of precisely this concept, and then believed that it was the lever,
is not difficult to guess. When Prof. Martensen
advanced with his great dogmatics, and wanted to make
room to right and left, the concept ``coherent
cognition'' was precisely the warning cry with which he signified to ``the absolutes in faith'' and ``the absolutes in knowledge''
that they now had to step aside, and not come too
near him with their ``stray thoughts and aphorisms, fancies
and flashes.'' When I now in this crush was seized
with fright, and was in anxiety at being pitifully
run over, I grabbed for the ``lever'', and fortunately found it
in ``the paradox'' in ``Johannes Climacus''. In order,
however, to remind the reader of the veiled challenge
that seemed to ring so bravely in the Martensenian slogan,
I let the concept
``coherent cognition'' run along in my ``examining review'' as a little
glaneur, so that in case the ``lever'' had
its effect, no ear of the speculative load should be wasted.
How, for the rest, I imagine I
have used the ``lever'' is clear enough from my last
thesis: ``\emph{Christian dogmatics is}
% --- p. 23 ---
\setcounter{footnote}{0}%
\opage{23}\emph{without principle}''. --- That Prof. Scharling has not hit upon
seeking the ``lever'' in this thesis certainly seems most
strange, since the bar is otherwise palpable enough; but, as
said, I greatly fear that His Rev. has his own opinions
about what a principle is.

That I could in no way use the concept
``coherent cognition'' as a ``lever'',
anyone will see who knows at all what a lever is.
A lever must after all be firm and strong; but the concept ``coherent
cognition'' is precisely a very soft and pliable
concept that winds itself through all possible degrees. There
are after all many kinds of coherence and many degrees
in every coherence; there is real coherence,
ideal coherence, firm coherence, loose coherence:
everything, even the genial Apology, has its coherence; if Prof. Scharling therefore thinks he has found a bar in the category
coherence, it can
certainly not be a ``lever''; it must then presumably be --- another sort of bar; yet surely never a lime-rod!

As reward for his noble striving to boss ``the paradox'' about
and keep order in literature, Prof. Scharling expects
that ``the men against whose conduct he has now declared
himself'' will cool ``their indignation in a philippic ---
whereby he foresees a treatment no less harsh
than that which Marsyas experienced from the angered god.'' ---
Alas! how strangely things go in this world
with human foresight. I foresaw quite surely
that I would get an answer from Prof. Martensen,
and see, then I get --- an Apology from Prof. Scharling;
Prof. Scharling in turn foresees that he will
get ``a philippic'' from me, and see, then he gets
--- a few genial words.

It is, by the way, strange enough that Prof. S. has
been able to ``foresee'' that his genial Apology should really
% --- p. 24 ---
\setcounter{footnote}{0}%
\opage{24}deserve a philippic, and thereby also foresee that it would
go with him as with Marsyas; this foresight must then surely also
stem from one of the dark possibilities he sometimes
sees when he sees the paradoxes in Pandora's box.
As far as I am concerned, His Rev. can truly rest assured
that I shall never visit him with any philippic,
or cool my indignation on an innocent man who
surely will never arouse my indignation. In my writings I have
never expressed any indignation against Prof. Scharling, never
challenged Prof. Scharling, and, to put it plainly, seldom % print: reentud (misprint)
thought of Prof. Scharling. If Prof. S. will nonetheless
kindly continue, as hitherto, to think of me, set
me a little to rights and take care of my upbringing,
I shall also know how to appreciate it in all quietness; only
Prof. S. must not take it amiss if he only seldom sees
written proof of my gratitude; I have, as
said, a fear of engaging too much with Prof.
Scharling, who always alarms me with his dark possibilities.


That the possibilities may be dark enough,
I, who am not safe from any side, but only subsist
by believing the best and fearing the worst, should
least of all presume to deny. But even
if the worst should happen, even if it should one day
happen that either ``the paradox'' or ``speculation'' found
itself entitled to endorse my servant's testimonial, dismiss
me as a useless servant, and pack me off ---
``nach Hause''; then I would nevertheless, as I am packed off ``nach
Hause'', quite gladly take along the consolation that ``theological
science'' must at last have found its principle, and
therewith also have preserved the conviction that Professor
Scharling's ``spiritual nature constitutes a unity.''

\begin{center}\rule{4em}{0.4pt}\end{center}

\end{document}
